% explicitness.tex - Explicit capability specifications

\section{Should Capabilities Be Explicit?}
\label{sec:explicit-capabilities}

A central design question is how visible capabilities should be. %
One philosophy is that capabilities should remain visible in the program
structure, since visible authority supports review, audits, and a clear
security story. %
In private communication, the security expert Matt Rice has raised the concern
that hiding capability propagation can make it harder for reviewers to see
where authority flows. %

\lambdaCC{} leaves this choice to language and library designers. %
Because capabilities are statically tracked, a language can make capability
requirements explicit where that visibility matters, while still avoiding
manual propagation through every intermediate function. %
For example, a language may require public functions to have explicit capability
specifications. %
Such usability concerns are up to language designers' discretion. %
We believe it is good practice to make capability requirements explicit at
module boundaries. %
Within a module, contextual binding can reduce repetitive plumbing without
turning capabilities into ambient authority. %
